\section{Marco legal y normativo}
\label{sec:marco-legal}

La operación de un sistema que captura tráfico de red y registra 
información sobre los dispositivos conectados debe ajustarse al marco 
legal aplicable a la protección de los datos personales y al acceso 
a sistemas informáticos. Este ajuste se evalúa tradicionalmente desde 
los tres pilares de la seguridad de la información, conocidos como 
tríada CIA: confidencialidad, que exige que la información solo sea 
accesible a quienes tienen autorización explícita; integridad, que 
exige que los datos no sean alterados de forma indebida durante su 
tratamiento; y disponibilidad, que exige que la información se 
mantenga accesible para los usuarios legítimos en el momento en que 
la requieran. Las decisiones de diseño descritas en la Sección 
\ref{sec:desarrollo}, como la sanitización de payload, la 
autenticación basada en tokens firmados, el cifrado en tránsito y la 
separación de procesos para aislar fallos, responden de forma directa 
a estos tres pilares.

En el ordenamiento jurídico colombiano, las normas aplicables a 
GOATGuard se concentran en tres bloques: la tipificación de los 
delitos contra la confidencialidad, integridad y disponibilidad de 
los datos y los sistemas informáticos \cite{ley1273_2009}; el régimen 
general de protección de datos personales \cite{ley1581_2012} y su 
decreto reglamentario parcial \cite{decreto1377_2013}; y las 
disposiciones sobre monitoreo laboral \cite{ley2466_2025}. La 
Tabla~\ref{tab:marco-legal} sintetiza los artículos relevantes y la 
forma en que el diseño del sistema responde a cada uno de ellos.

\begin{table*}[!t]
\centering
\caption{Normativa colombiana aplicable al sistema y su impacto en el diseño.}
\label{tab:marco-legal}
\begin{tabular}{@{}p{2.5cm}p{1.8cm}p{5cm}p{6cm}@{}}
\toprule
\textbf{Norma} & \textbf{Artículo} & \textbf{Materia regulada} & \textbf{Respuesta de diseño en GOATGuard} \\
\midrule
Ley 1273 de 2009 & 269A & Acceso abusivo a sistema informático & La instalación del sensor requiere autorización expresa del responsable de la red. \\[4pt]
Ley 1273 de 2009 & 269C & Interceptación de datos informáticos & La captura de paquetes se realiza bajo autorización documentada; el payload se sanitiza para no interceptar contenido de comunicaciones. \\[4pt]
Ley 1581 de 2012 & 4 (libertad) & Consentimiento del titular de los datos & Los usuarios de los endpoints son informados del monitoreo antes de la puesta en operación. \\[4pt]
Ley 1581 de 2012 & 4 (finalidad) & Tratamiento únicamente para el fin declarado & Los datos capturados se usan solo para la detección de anomalías y el inventario de activos, no para otros propósitos. \\[4pt]
Ley 1581 de 2012 & 6 & Excepción para fines académicos o de investigación & Aplicable al presente proyecto, siempre que se mantenga la anonimización del contenido. \\[4pt]
Decreto 1377 de 2013 & 3 & Aviso previo de privacidad & El despliegue del sistema incluye la entrega de un aviso de privacidad a los usuarios antes de iniciar la captura. \\[4pt]
Ley 2466 de 2025 & N/A & Monitoreo en el ámbito laboral & En escenarios empresariales, la operación del sistema requiere el consentimiento explícito y documentado del trabajador. \\
\bottomrule
\end{tabular}
\end{table*}

Tres decisiones concretas del diseño resultan del cruce con la 
tabla anterior. La primera es la sanitización del payload en el 
sensor de captura: el recorte dinámico de los paquetes preserva los 
encabezados necesarios para el análisis de flujos pero elimina el 
contenido de las comunicaciones, lo que reduce el alcance de la 
captura al estrictamente necesario para el fin declarado y elimina la 
exposición a contenido sensible que podría constituir una 
interceptación indebida en el sentido del artículo 269C. La segunda 
es el control de acceso al sistema mediante autenticación con 
contraseña hasheada y emisión de tokens JWT con expiración definida, 
que limita el acceso a los datos procesados únicamente al 
administrador autorizado, en concordancia con el principio de 
confidencialidad. La tercera es la generación de un registro 
persistente de las operaciones realizadas por el pipeline y la API, 
que permite auditar el uso del sistema y respaldar la trazabilidad 
exigida por el régimen general de protección de datos.

La aplicación académica del sistema se encuadra dentro de la 
excepción científica del artículo 6 de la Ley 1581 de 2012, 
condicionada a la anonimización del tratamiento. El equipo adopta 
esta condición como criterio operativo: ningún identificador de 
usuario individual se persiste en el sistema, y los identificadores 
de dispositivos se construyen a partir de atributos técnicos 
(hostname y dirección MAC) que no permiten por sí mismos vincular la 
actividad observada con una persona natural específica. La extensión 
del sistema a entornos productivos, donde esta excepción no aplicaría 
de forma automática, requeriría la gestión explícita de los 
consentimientos correspondientes según la Ley 2466 de 2025 en 
contextos laborales.

La totalidad de las actividades de captura, procesamiento y pruebas 
ejecutadas durante el desarrollo del proyecto se circunscriben a 
entornos controlados dispuestos por el equipo de investigación, 
conformados por redes locales de uso académico y por los endpoints 
propios de los integrantes del proyecto. Los datos capturados durante 
estas pruebas no se transfieren a terceros, no se integran a 
sistemas en producción y no se emplean para ningún fin ajeno al de 
la validación técnica del prototipo. La eventual adopción del sistema 
en un entorno productivo o empresarial queda por fuera del alcance 
del presente trabajo y requeriría, como condición previa, la 
formalización de los consentimientos, autorizaciones y controles 
adicionales que el marco normativo señalado exige para el tratamiento 
de datos personales fuera de la excepción científica.